% =====================================================================
% Paper CONAIISI — pipeline de dataset y entrenamiento de pose
% ESQUELETO (Fase 0): títulos de sección + marcadores \pendiente{} en su
% lugar. NO tiene prosa todavía — se escribe en una pasada posterior con
% aprobación previa del usuario. Ver datos_extraidos.md (hechos con fuente
% verificada) y PENDIENTES.md (huecos agrupados por sección) en esta misma
% carpeta.
%
% Formato clonado de UADE_PFI_Template-develop/paper_conaiisi/paper_conaiisi.tex
% (revisión doble ciego CoNaIISI, A4 2 columnas, Times 10 pt, máx. 10 págs.
% incluidas referencias). Ese archivo es el "paper hermano" (comparador de
% movimiento / DTW): este paper es su complemento y toma esa capa como
% fuera de alcance.
%
% BUSCAR "\pendiente" antes de escribir prosa final o de enviar: son los
% huecos de contenido y los marcadores de cita pendiente de verificar.
% =====================================================================
\documentclass[10pt,a4paper,twocolumn]{article}

\usepackage[utf8]{inputenc}
\usepackage[T1]{fontenc}
\usepackage[spanish,es-nodecimaldot,es-tabla]{babel}
\usepackage{mathptmx}            % Times para texto y matemática
\usepackage{helvet}              % Helvética para pies de figura y tabla
\usepackage{courier}             % Courier para código fuente
\usepackage{amsmath}
% ----- Página: área de impresión y columnas ---------------------------
% Idéntico a paper_conaiisi.tex: columnas de 8,25 cm con 8 mm de separación.
\usepackage[a4paper,
  textwidth=17.3cm,     % 2 x 8,25 cm + 8 mm
  top=3cm,              % margen superior exacto
  textheight=22.54cm,   % alto del área de impresión
  columnsep=8mm]{geometry}
\pagestyle{empty}       % sin número de página: nada fuera del área de impresión
\usepackage{titlesec}
\usepackage{caption}
\usepackage{booktabs}
\usepackage{array}
\usepackage{graphicx}
\usepackage{tikz}
\usetikzlibrary{positioning,arrows.meta,fit}
\usepackage{xcolor}
\usepackage{url}

% --- Tipografía de títulos según plantilla -----------------------------
\titleformat{\section}{\normalfont\fontsize{12}{14}\selectfont\bfseries}{\thesection.}{0.5em}{}
\titlespacing*{\section}{0pt}{24pt}{12pt}
\titleformat{\subsection}{\normalfont\fontsize{11}{13}\selectfont\bfseries}{\thesubsection}{0.5em}{}
\titlespacing*{\subsection}{0pt}{6pt}{6pt}
\titleformat{\subsubsection}{\normalfont\fontsize{10}{12}\selectfont\bfseries}{\thesubsubsection}{0.5em}{}
\titlespacing*{\subsubsection}{0pt}{12pt}{6pt}

% ----- Notas al pie: Times 9 pt, interlineado simple -------------------
\renewcommand{\footnotesize}{\fontsize{9}{11}\selectfont}

% --- Pies de figura y títulos de tabla según plantilla oficial ---------
\DeclareCaptionFont{nueve}{\fontsize{9}{11}\selectfont}
\captionsetup{font={sf,bf,nueve},labelsep=period,justification=centering}
\captionsetup[figure]{position=below,skip=4pt}
\captionsetup[table]{position=above,skip=4pt}

% --- Párrafos ----------------------------------------------------------
\setlength{\parindent}{5mm}
\setlength{\parskip}{0pt}

% --- Marcador de contenido pendiente (quitar antes del envío final) ----
\newcommand{\pendiente}[1]{\textcolor{red}{\textbf{[PENDIENTE: #1]}}}
\newcommand{\verificarcita}[1]{\textcolor{red}{\textbf{[VERIFICAR CITA: #1]}}}

\begin{document}

% ======================= TÍTULO Y AUTORES ==============================
% Revisión doble ciego: versión anónima. Para la versión final, comentar
% las líneas "anonimizado" y descomentar el bloque de autores (mismo
% formato/autores que paper_conaiisi.tex: Fernández, Witenas, Martínez,
% UADE Facultad de Ingeniería).
\twocolumn[{%
\centering
\vspace*{0.49cm}
{\fontsize{14}{17}\selectfont\bfseries
% TÍTULO TENTATIVO — confirmar con el resto del equipo antes de la versión final.
Construcción de un Conjunto de Datos Densamente Anotado y Compresión de un Estimador de Pose para Ejercicios de Rehabilitación en Dispositivos Móviles\par}
\vspace{2\baselineskip}
{\fontsize{12}{14}\selectfont Autores anonimizados para revisión\par}
{\fontsize{12}{14}\selectfont\itshape Filiación anonimizada para revisión\par}
% ---- BLOQUE DE AUTORES PARA LA VERSIÓN FINAL (descomentar) ------------
% \noindent
% \begin{minipage}[t]{0.32\textwidth}\centering
% {\fontsize{12}{14}\selectfont Fernández,\\Nicolás Esteban\par}
% \vspace{2pt}
% {\fontsize{12}{14}\selectfont\itshape Universidad Argentina de\\la Empresa (UADE),\\Facultad de Ingeniería\par}
% \vspace{2pt}
% {\fontsize{10}{12}\selectfont nicolasefernandez@uade.edu.ar\par}
% \end{minipage}%
% \begin{minipage}[t]{0.32\textwidth}\centering
% {\fontsize{12}{14}\selectfont Witenas,\\Juan Ignacio\par}
% \vspace{2pt}
% {\fontsize{12}{14}\selectfont\itshape Universidad Argentina de\\la Empresa (UADE),\\Facultad de Ingeniería\par}
% \vspace{2pt}
% {\fontsize{10}{12}\selectfont jwitenas@uade.edu.ar\par}
% \end{minipage}%
% \begin{minipage}[t]{0.32\textwidth}\centering
% {\fontsize{12}{14}\selectfont Martínez,\\María Roxana\par}
% \vspace{2pt}
% {\fontsize{12}{14}\selectfont\itshape Universidad Argentina de\\la Empresa (UADE),\\Facultad de Ingeniería\par}
% \vspace{2pt}
% {\fontsize{10}{12}\selectfont mariarmartinez@uade.edu.ar\par}
% \end{minipage}\par
% -----------------------------------------------------------------------
\vspace{2\baselineskip}
}]

% ======================= RESUMEN =======================================
\begin{center}
{\fontsize{12}{14}\selectfont\bfseries Resumen}
\end{center}
\begingroup\itshape
\noindent Este trabajo presenta la construcción del conjunto de datos y el
entrenamiento de un estimador de pose comprimido para ejercicios de
rehabilitación, destinado a ejecutarse íntegramente en un teléfono Android de
gama media. El dominio objetivo ---adultos mayores, ejercicios en
sedestación, cámara frontal en el hogar--- no está cubierto por los corpus
públicos de pose humana. Se adoptó una estrategia de anotación en dos
niveles: pseudo-etiquetado denso sobre la totalidad del corpus con un
anotador de alta capacidad, y una regla de derivación a revisión humana
dirigida solo a la cola de casos de baja confianza. El modelo se comprime
mediante poda estructurada, recuperación, destilación de conocimiento y
cuantización entera, distinguiendo el artefacto desplegable del instrumento
---la validación cruzada agrupada por sujeto--- que mide su generalización.
Se reporta la composición del corpus, la configuración de la cadena de
compresión y el diagnóstico de dos hallazgos negativos de la primera
iteración: un colapso de exactitud fuera de dominio pese a un desempeño
saludable en la partición de validación, y una etapa de destilación cuya
degradación se rastreó hasta tres defectos concretos de implementación.
\endgroup

\vspace{6pt}
\noindent{\fontsize{10}{12}\selectfont\bfseries Palabras clave:} anotación asistida
por modelo, pseudo-etiquetado, estimación de pose humana, compresión de modelos,
aprendizaje profundo en el borde, rehabilitación domiciliaria.

% ======================= 1. INTRODUCCIÓN ===============================
\section{Introducción}

La adherencia a un programa de rehabilitación motriz depende en buena medida
de lo que ocurre fuera del consultorio, en las sesiones domiciliarias que el
profesional no supervisa en tiempo real. En personas adultas mayores esa
ausencia de supervisión se traduce en errores de ejecución que pasan
inadvertidos y en abandono de la rutina, y la telerehabilitación asistida por
video ha mostrado evidencia de efectividad frente a ese problema
\cite{jirasakulsuk2022}. Buena parte de esa evidencia depende, sin embargo,
de que el propio sistema pueda decidir si una ejecución es correcta, lo que
exige antes que nada un estimador de pose confiable corriendo sobre el
dispositivo del paciente. En un trabajo previo \cite{anonimo_previo} se
presentó una plataforma móvil Android orientada a este escenario, que procesa
el video de la cámara frontal del teléfono en tiempo real y sin conexión a
internet; ese trabajo describió la arquitectura general del sistema, pero
dejó sin desarrollar la construcción del propio modelo de estimación de pose
que la sostiene.

Este artículo se ocupa precisamente de esa construcción: el conjunto de
datos, la estrategia de anotación y el proceso de entrenamiento y compresión
que producen el estimador de pose que la aplicación ejecuta. El componente
que consume esa salida para decidir si una ejecución del paciente es
correcta ---basado en alineamiento temporal de trayectorias angulares--- se
desarrolla en un trabajo complementario \cite{anonimo_dtw} y queda fuera del
alcance de este artículo, que toma esa capa posterior como dada.

Las contribuciones de este trabajo son seis. Primero, la caracterización del
dominio objetivo ---adultos mayores, ejercicios en sedestación, cámara
frontal de teléfono en interiores domésticos--- y de su brecha frente a los
corpus generales de pose, con la restricción de uso que ese dominio impone
sobre el diseño del pipeline de anotación, enunciada como una decisión de
diseño y no como un hecho de licencia verificado (\S\ref{sec:requisitos}).
Segundo, una estrategia de anotación en dos niveles: pseudo-etiquetado masivo
con un criterio de calibración de umbrales medido sobre una muestra
estratificada, y un criterio explícito de derivación a revisión humana que
concentra el esfuerzo manual en la cola de casos difíciles en lugar de
distribuirlo sobre todo el corpus. Tercero, la distinción operativa entre el
modelo que efectivamente se despliega ---entrenado una única vez sobre una
partición fija--- y el instrumento que mide su capacidad de generalizar a
personas no vistas ---la validación cruzada agrupada por sujeto---, evitando
que ambos se confundan en el registro de experimentos. Cuarto, una cadena de
compresión con un orden justificado (poda, recuperación, destilación,
cuantización) y un contrato explícito para el artefacto final. Quinto, el
diagnóstico honesto de dos hallazgos negativos de la primera iteración del
modelo: un colapso de exactitud fuera de dominio pese a un desempeño
saludable en la partición de validación, y una etapa de destilación cuya
degradación se rastreó hasta defectos concretos de implementación. Sexto, un
conjunto de prácticas de reproducibilidad ---versionado de datos, registro de
experimentos, determinismo--- que atraviesan todo el proceso.

El resto del artículo se organiza así: la sección 2 revisa los trabajos
relacionados; la sección 3 desarrolla los requisitos del conjunto de datos;
la sección 4 describe el pipeline de anotación; la sección 5 presenta el
entrenamiento y la cadena de compresión; la sección 6 reporta los resultados,
incluido el diagnóstico de la iteración fallida; la sección 7 discute las
limitaciones y la sección 8 concluye.

% ======================= 2. TRABAJOS RELACIONADOS ======================
\section{Trabajos relacionados}

\subsection{Datasets de pose en rehabilitación y sus limitaciones}
Los corpus de referencia en estimación de pose, como COCO \cite{lin2014}, se
componen mayormente de escenas cotidianas o deportivas con personas de pie,
de cuerpo casi completo y sobre una amplia variedad de fondos; ninguno de
esos supuestos se sostiene en rehabilitación domiciliaria, donde el sujeto
está sentado, parcialmente ocluido por el mobiliario y filmado de frente por
la cámara de un teléfono. El dataset UCO Physical Rehabilitation
\verificarcita{atribución del dataset UCO Physical Rehabilitation — no
confirmada, ver PENDIENTES.md \S2} es, hasta donde se relevó para este
trabajo, uno de los pocos que anota ejercicios de rehabilitación cuadro a
cuadro sobre video; su limitación principal, que motiva la sección
\ref{sec:pipeline}, es que anota únicamente las articulaciones relevantes a
cada ejercicio ---típicamente tres de los diecisiete puntos del esquema
COCO--- y dejar el resto sin etiquetar impide entrenar directamente un
estimador de cuerpo completo.

\subsection{Anotación asistida y pseudo-etiquetado}
Cuando el costo de anotar manualmente un corpus completo es prohibitivo, una
alternativa establecida es generar pseudo-etiquetas con un modelo ya
entrenado y usarlas como señal de entrenamiento de otro \cite{lee2013}. La
variante adoptada en este trabajo es un anotador \emph{top-down} de dos
etapas ---un detector de personas seguido de un estimador de pose sobre el
recorte--- \verificarcita{"top-down es el enfoque de mayor exactitud
publicado en estimación de pose" — no hay una fuente concreta identificada,
ver PENDIENTES.md \S2}, instanciado con un detector de tiempo real
\cite{lyu2022} y un estimador basado en Vision Transformers \cite{xu2022},
ejecutados sobre el framework de OpenMMLab \cite{mmpose2020}. A diferencia
del uso más habitual del pseudo-etiquetado como forma de destilación de
conocimiento hacia un modelo más chico, aquí el anotador nunca se despliega:
su única función es producir etiquetas, y su costo computacional es
irrelevante porque corre una única vez fuera de línea.

\subsection{Compresión de modelos para dispositivos móviles}
La poda estructurada elimina canales completos de la red siguiendo su grafo
de dependencias \cite{fang2023}, sobre la idea original de podar y luego
recuperar precisión con reentrenamiento \cite{han2015}. La destilación de
conocimiento transfiere la capacidad de un modelo grande a uno chico
comparando sus salidas \cite{hinton2015} o, en la variante de transferencia
de atención adoptada en este trabajo, sus mapas de activación intermedios
\cite{zagoruyko2017}. La cuantización entera post-entrenamiento reduce la
precisión numérica de pesos y activaciones con una degradación de exactitud
acotada \cite{jacob2018}. Estas tres técnicas se combinan habitualmente para
producir un artefacto final compatible con hardware móvil, pero la
literatura rara vez discute el orden en que deben aplicarse cuando la
recuperación de precisión no puede apoyarse en un dataset completamente
etiquetado ---la situación exacta de este trabajo (\S\ref{sec:entrenamiento}).

\subsection{Métricas de evaluación de pose}
La similitud de keypoint objeto (OKS) y el porcentaje de puntos correctos
(PCK) sobre el esquema de diecisiete puntos de COCO \cite{lin2014} son las
métricas estándar para reportar exactitud de un estimador de pose, y ambas se
usan en este trabajo por etapa de la cadena de compresión. Un aspecto menos
estandarizado es cómo particionar los datos para medir generalización: un
particionado aleatorio a nivel de cuadro sobreestima el desempeño cuando el
mismo sujeto aparece en ambos lados de la partición, un caso particular del
problema más general de aproximar el caso de uso real en aprendizaje
automático clínico \cite{saeb2017}; la validación cruzada agrupada por
sujeto (\S\ref{sec:entrenamiento}) es la respuesta adoptada.

% ======================= 3. REQUISITOS DEL DATASET ======================
\section{Requisitos del conjunto de datos}
\label{sec:requisitos}

\subsection{Caracterización del dominio objetivo}
El sistema descripto en \cite{anonimo_previo} está pensado para adultos
mayores que ejecutan ejercicios de rehabilitación sentados, en su domicilio,
filmados con la cámara frontal de un teléfono apoyado a corta distancia. Esa
combinación de condiciones —sedestación, oclusión parcial por el respaldo de
una silla o por la ropa holgada habitual en adultos mayores, encuadre cercano
y de frente, iluminación doméstica no controlada— es exactamente la que los
corpus generales de pose humana no cubren: esos corpus están dominados por
personas de pie, de cuerpo completo, en escenas exteriores o deportivas. Un
estimador entrenado solo sobre esa distribución no tiene garantías de
generalizar al dominio objetivo, y verificarlo empíricamente es, de hecho, uno
de los hallazgos de este trabajo (\S\ref{sec:resultados}).

\subsection{Restricciones de uso y su traducción a requisitos técnicos}
El diseño del pipeline de anotación asume, como restricción de partida, que
los datos del dominio de rehabilitación no pueden subirse a servicios de
anotación o de entrenamiento de terceros. \pendiente{Verificar los términos
exactos de la licencia del dataset UCO Physical Rehabilitation contra su sitio
de distribución oficial o su publicación original antes de presentar esta
restricción como un hecho legal; hasta entonces se la trata como una decisión
de diseño asumida por el equipo, no como una cláusula contractual confirmada.}
Bajo esa restricción, tres requisitos técnicos se vuelven no negociables: la
anotación asistida debe ejecutarse sobre infraestructura propia y no sobre una
plataforma de anotación alojada por terceros, lo que motiva el uso de una
instancia autoalojada de una herramienta de anotación de código abierto
(\S\ref{sec:pipeline}); el almacenamiento de datos y modelos debe quedar
versionado en un remoto controlado por el propio equipo y no en un servicio
de entrenamiento externo; y ningún artefacto que se publique como resultado
de este trabajo —el modelo exportado, las métricas— puede permitir
reconstruir imágenes o identidades de los sujetos filmados.

\subsection{Criterios de selección del dataset}
La selección de UCO Physical Rehabilitation como base para este trabajo
respondió a cinco criterios. El primero y excluyente es la correspondencia de
dominio: se necesitan ejercicios de rehabilitación, no escenas genéricas. El
segundo es la anotación por cuadro sobre video —y no sobre imágenes
sueltas—, porque el caso de uso es inferencia continua sobre una secuencia de
movimiento. El tercero es el volumen: la cadena de compresión de la sección
\ref{sec:entrenamiento} necesita decenas de miles de cuadros para el ajuste
fino y la destilación. El cuarto es que el dataset identifique explícitamente
a qué sujeto pertenece cada cuadro, condición sin la cual el particionado por
sujeto de la sección \ref{sec:entrenamiento} no es posible. El quinto es la
disponibilidad efectiva para investigación, tanto práctica como legal —esta
última sujeta a la verificación pendiente de la subsección anterior—. El
dataset satisface los primeros cuatro sin ambigüedad; su limitación conocida
de antemano, que la anotación cubre solo entre dos y tres articulaciones por
ejercicio, es precisamente el problema que resuelve la sección
\ref{sec:pipeline}.

% ======================= 4. PIPELINE DE ANOTACIÓN =======================
\section{Pipeline de anotación}
\label{sec:pipeline}

La Figura~\ref{fig:pipeline} resume el recorrido de los datos, desde el pool
crudo hasta los conjuntos de entrenamiento particionados por sujeto.

\begin{figure*}[t]
\centering
\resizebox{\textwidth}{!}{%
\begin{tikzpicture}[
  font=\footnotesize,
  caja/.style={draw, rounded corners=2pt, align=center, minimum height=8mm, inner sep=3pt},
  cajapend/.style={draw, rounded corners=2pt, align=center, minimum height=8mm, inner sep=3pt, dashed, text=black!55},
  etiqueta/.style={font=\scriptsize\itshape, text=black!60},
  flecha/.style={-{Stealth[length=2mm]}, thick},
  flechapend/.style={-{Stealth[length=2mm]}, thick, dashed},
  node distance=4mm and 6mm
]
% --- fila principal ---
\node[caja] (raw) {Pool UCO crudo\\(27 sujetos, 10 ejerc.,\\5 cámaras, 1280$\times$720)};
\node[caja, right=of raw] (conv) {Conversor\\(stride=30, bbox\\hull+5\%, 74.107 pares)};
\node[caja, right=of conv] (pseudo) {Pseudo-etiquetado\\RTMDet-m + ViTPose-H};
\node[caja, right=of pseudo] (manifest) {Manifiesto\\(flags de confianza)};
\node[caja, right=of manifest] (kfold) {GroupKFold\\por sujeto (k=5)};
\node[caja, right=of kfold] (train) {Entrenamiento\\y compresión (\S\ref{sec:entrenamiento})};
% --- fila inferior: cola de auditoría manual ---
\node[caja, below=11mm of pseudo] (baja) {Cuadros de\\baja confianza};
\node[caja, right=of baja] (cvat) {Auditoría CVAT\\(selección + corrección\\humana + round-trip gate)};
\node[cajapend, right=of cvat] (gold) {Conjunto dorado\\\pendiente{auditoría no ejecutada aún}};
% --- flechas fila principal ---
\draw[flecha] (raw) -- (conv);
\draw[flecha] (conv) -- (pseudo);
\draw[flecha] (pseudo) -- (manifest);
\draw[flecha] (manifest) -- (kfold);
\draw[flecha] (kfold) -- (train);
% --- flechas cola de auditoría ---
\draw[flecha] (manifest.south) to[out=-90,in=90] (baja.north);
\draw[flecha] (baja) -- (cvat);
\draw[flechapend] (cvat) -- (gold);
% --- etiquetas de fila ---
\node[etiqueta, above=1mm of raw.north west, anchor=south west] {Curado y etiquetado automático};
\node[etiqueta, below=1mm of baja.south west, anchor=north west] {Cola de revisión humana (dirigida por confianza)};
\end{tikzpicture}%
}
\caption{Pipeline de construcción del conjunto de datos. Los cuadros de baja
confianza se derivan a auditoría manual en CVAT; el conjunto dorado resultante
(recuadro punteado) todavía no existe porque la auditoría no se ejecutó (ver
\S\ref{sec:resultados}). Elaboración propia.}
\label{fig:pipeline}
\end{figure*}

\subsection{Curado y muestreo}
La conversión del material crudo a un formato de entrenamiento estandarizado
integra en una sola pasada la extracción de cuadros y la generación de
etiquetas, garantizando que cada imagen tenga su etiqueta correspondiente sin
desfasajes. El muestreo temporal retiene un cuadro cada treinta, porque
cuadros consecutivos de un movimiento lento son casi idénticos y retenerlos
todos multiplicaría el costo de almacenamiento sin aportar variedad; con ese
paso, el conjunto en alcance —veintisiete sujetos ejecutando ejercicios de
flexo-extensión de rodilla y pierna y de elevación de brazo, registrados por
cinco cámaras— produce 74.107 pares de imagen y etiqueta. Cada cuadro
retenido se identifica con un esquema determinístico que codifica sujeto,
ejercicio, cámara y número de cuadro, con una función de análisis inversa que
rechaza explícitamente cualquier identificador malformado; ese identificador
es lo que luego permite recuperar el sujeto para el particionado de la
sección \ref{sec:entrenamiento} sin tener que inferirlo de otra fuente. Las
doce articulaciones que UCO anota se proyectan una a una sobre los puntos 5 a
16 del esquema COCO de diecisiete puntos, mientras que los cinco puntos
faciales —que UCO nunca anota— se escriben como no visibles; la caja
delimitadora de cada persona se deriva de la envolvente de los puntos
visibles con un margen del 5\,\%.

\subsection{Pseudo-etiquetado masivo}
Etiquetar manualmente las 74.107 imágenes del pool con los diecisiete puntos
completos es impracticable a la escala de este trabajo, y entrenar
únicamente con la anotación esparsa original —dos o tres puntos por
cuadro— enseña al modelo a declarar ausente el resto del cuerpo, que sí está
visible. La alternativa adoptada es generar pseudo-etiquetas densas sobre la
totalidad del pool con un anotador \emph{top-down} de dos etapas
\cite{lee2013}: un detector de personas en tiempo real \cite{lyu2022}
localiza al sujeto y un estimador de pose basado en Vision Transformers
\cite{xu2022} predice los diecisiete puntos sobre el recorte, ambos ejecutados
con el stack de OpenMMLab \cite{mmpose2020}. Los dos umbrales de confianza del
anotador se calibraron en el orden que su dependencia impone y no al azar:
primero el umbral de detección, porque una detección perdida descarta el
cuadro entero y es irrecuperable, y solo después el umbral de visibilidad de
keypoints, que decide qué entra a la función de pérdida del modelo
estudiante. Un barrido sobre una muestra de 500 imágenes estratificada por
sujeto y ejercicio mostró que un umbral de detección de 0,3 pierde apenas el
0,2\,\% de los cuadros —frente al 1,0\,\% de un umbral de 0,5—, sin admitir
detecciones espurias porque estas quedan igualmente sujetas al filtro
posterior de keypoints; sobre ese punto de operación, el percentil 5 de
confianza de keypoints nunca cayó por debajo de 0,41 (los peores casos, codos
y muñecas), mientras que el miembro inferior —crítico para los ejercicios en
alcance— se mantuvo entre 0,75 y 0,82. La corrida completa, ejecutada una
única vez sobre una GPU RTX 3090 en aproximadamente 2,1 horas, escribió
73.955 personas sobre las 74.107 imágenes, de las cuales el 94,3\,\%
(69.880) quedaron con los diecisiete puntos completos y solo 152 imágenes no
tuvieron ninguna detección, con una confianza media de keypoint de 0,8897
sobre todo el pool.

\subsection{Criterio de derivación a revisión humana}
El pseudo-etiquetado automático no es perfecto, y el aporte de este trabajo en
ese punto no es una corrección manual exhaustiva sino un criterio explícito
de dónde concentrarla. Cada imagen queda marcada como de baja confianza si su
confianza media de keypoint cae por debajo de 0,5 o si retiene menos de diez
de los diecisiete puntos visibles, y se registra el motivo específico —sin
detección, sin keypoints válidos, múltiples personas en el cuadro o
confianza baja con una sola persona—. Ese subconjunto, más una muestra de
control de 150 cuadros sanos estratificada por sujeto y ejercicio (y, dentro
de cada ejercicio, por si se trata de un movimiento en sedestación o de
elevación de brazo), define el alcance completo de la revisión humana: en
lugar de auditar una muestra aleatoria del pool, donde la enorme mayoría de
las etiquetas ya es correcta y el esfuerzo se diluye, la revisión se
concentra exactamente donde el propio anotador declara incertidumbre. El
resultado cuantitativo de aplicar ese criterio —cuántos cuadros se
corrigieron y con qué costo— se reporta en la sección \ref{sec:resultados}.

\subsection{Cola de anotación manual}
La infraestructura de revisión humana corre sobre una instancia propia de una
herramienta de anotación de código abierto, en línea con la restricción de no
depender de servicios de terceros (\S\ref{sec:requisitos}). Antes de admitir
cualquier corrección real, un control de ida y vuelta —importar un lote
reducido de cuadros, exportarlo sin tocar nada y comparar el resultado contra
el original— verificó que la herramienta preserva la visibilidad de cada
punto y que el ruido de fondo introducido por el propio ida y vuelta es de
apenas 0,005 píxeles, muy por debajo del umbral que separa una corrección
real de un artefacto de la herramienta. Una corrección se define como un
desplazamiento de un punto por encima de ese umbral o un cambio de su
visibilidad entre ausente y presente, y el resultado de la reconciliación
produce, además de una métrica de tasa de corrección por motivo y por tipo de
ejercicio, un conjunto dorado de imágenes verificadas que queda excluido del
entrenamiento y reservado para evaluación. \pendiente{La ejecución real de la
auditoría —selección del subconjunto, empaquetado, corrección humana en la
herramienta y reconciliación— no se completó todavía sobre el corpus
completo; no existe aún el conjunto dorado ni el reporte de auditoría. Ver
PENDIENTES.md \S4 para el comando exacto que los generaría.}

% ======================= 5. ENTRENAMIENTO Y COMPRESIÓN ==================
\section{Entrenamiento y compresión para despliegue}
\label{sec:entrenamiento}

\subsection{Particionado por sujeto}
Un particionado aleatorio a nivel de cuadro es metodológicamente inválido
sobre este corpus: cuadros consecutivos de un mismo clip son casi idénticos y
el mismo sujeto aparece simultáneamente en cinco cámaras, de modo que un
reparto aleatorio dejaría cuadros casi duplicados a ambos lados de la
partición y las métricas resultantes medirían memorización de sujetos
vistos y no generalización a personas nuevas —exactamente la aproximación al
caso de uso real que la literatura de aprendizaje automático clínico señala
como necesaria \cite{saeb2017}—. El particionado adoptado es una validación
cruzada agrupada por sujeto con cinco grupos y semilla fija, de modo que cada
sujeto aporta cuadros a un único grupo en cada partición. El dataset
pseudo-etiquetado, generado de forma independiente al de etiquetas
originales, verifica en tiempo de generación que su reparto de sujetos por
grupo coincide exactamente con el de este último, y aborta ante cualquier
divergencia; sin esa verificación, la comparabilidad partición a partición
entre ambos conjuntos sería una suposición y no un hecho constatado.

\subsection{Configuración de entrenamiento}
El modelo estudiante es la variante más liviana de una arquitectura de una
sola pasada para detección y estimación conjunta de pose
\cite{maji2022,ultralytics2024}, entrenada a una resolución de entrada de 640
píxeles con optimizador AdamW y tasa de aprendizaje inicial de 0,01. Ambos
valores se declaran de forma explícita en la configuración de cada corrida
porque, en su modo por defecto, el framework de entrenamiento selecciona
automáticamente el optimizador y sobrescribe la tasa de aprendizaje recibida
por una derivada del tamaño del dataset —lo que habría descartado en
silencio la reducción deliberada de la tasa a un décimo que exige el ajuste
fino posterior a la poda (\S\ref{sec:cadena}).

\subsection{Modelo desplegable vs. instrumento de medición}
La validación cruzada de cinco particiones cumple una función de medición:
estima cuánto generaliza cada configuración a sujetos que el modelo nunca
vio, y es la que sostiene las comparaciones entre fases de este trabajo. El
modelo que efectivamente se integra a la aplicación, en cambio, se entrena
una única vez sobre una partición fija, con una configuración idéntica en
todo lo demás, y se registra en un experimento separado que deliberadamente
no reporta media ni desvío estándar entre particiones: una media calculada
sobre una única partición, con desvío cero, sería indistinguible en una tabla
de resultados de un resultado de validación cruzada saludable, y esa
confusión es exactamente la que la separación de experimentos busca evitar.
La distinción importa porque son dos artefactos con propósitos distintos
—uno mide, el otro se despliega— y no debería tratárselos como si fueran
intercambiables.

\subsection{Cadena de compresión}
\label{sec:cadena}
La compresión para el dispositivo aplica tres técnicas en un orden que no es
arbitrario. Primero, la poda estructurada elimina canales completos de la red
siguiendo su grafo de dependencias \cite{fang2023}, reduciendo los parámetros
del estudiante de 3,30 a 1,91 millones (un 42\,\% menos); medida sobre la
partición de validación completa, esa poda por sí sola lleva la exactitud
(OKS) de 0,846 a 0,000 —el modelo deja de detectar a la persona—, un colapso
total y no una degradación leve. Segundo, un ajuste fino de recuperación con
la tasa de aprendizaje reducida a un décimo de la original revierte ese
colapso. Tercero, la destilación de conocimiento transfiere al estudiante
podado y recuperado la capacidad de un modelo maestro de mayor tamaño
mediante transferencia de atención \cite{zagoruyko2017} —comparando los
mapas de activación intermedios de ambos en lugar de sus salidas finales
\cite{hinton2015}—, elegida en lugar de un ajuste fino supervisado clásico
por una razón específica de este trabajo: la transferencia de atención se
calcula únicamente a partir de imágenes y del modelo maestro, sin requerir
etiquetas ground-truth, mientras que las etiquetas parciales originales de
UCO —dos o tres de diecisiete puntos por ejercicio— enseñarían activamente
al modelo a suprimir el resto del cuerpo si se usaran para un ajuste fino
supervisado de cuerpo completo. La poda se aplica antes que cualquier
entrenamiento, de modo que tanto la recuperación como la destilación entrenan
ya la arquitectura reducida que llegará al dispositivo, y no una arquitectura
que se mutila después. Cuarto y último, la cuantización entera
post-entrenamiento \cite{jacob2018} convierte los pesos a enteros de 8 bits;
el artefacto que efectivamente se exporta cuantiza los pesos pero mantiene
las activaciones y la entrada/salida en punto flotante de 32 bits —una
variante de rango dinámico y no de enteros completos—, lo que en la práctica
es exactamente el contrato que espera la aplicación móvil.

\subsection{Infraestructura de cómputo}
Todo el entrenamiento y la cuantización se ejecutaron sobre una única GPU
NVIDIA RTX 3090 de 24 GB. Con el estudiante a resolución 640 y lote de 24, el
entrenamiento procesa aproximadamente 233 imágenes por segundo, es decir,
unos 4,1 minutos por época sobre los 57.805 cuadros de entrenamiento de una
partición. La asimetría más significativa está en la evaluación: un modelo
cuantizado evaluado sobre CPU —con lote fijo en una imagen, restricción
propia del formato de exportación— tarda 196 milisegundos por imagen, frente
a los 12,8 milisegundos de la misma arquitectura sin cuantizar sobre GPU; esa
diferencia es la que motiva acotar el conjunto de calibración de la
cuantización a un máximo de unas pocas centenas de imágenes representativas,
en lugar de recorrer con él la partición de validación completa.

% ======================= 6. RESULTADOS ===================================
\section{Resultados}
\label{sec:resultados}

\subsection{Estadísticas del corpus por etapa}
La Tabla~\ref{tab:corpus} resume la composición del conjunto de datos en cada
etapa del pipeline. De las 74.107 imágenes del pool, el pseudo-etiquetado
escribió una persona en 73.955 y dejó los diecisiete puntos completos en el
94,3\,\% de los casos, lo que deja una fracción acotada de cuadros —los
marcados de baja confianza más la muestra de control— que requieren atención
humana antes de poder usarse con la misma confianza que el resto del corpus.

\begin{table}[t]
\caption{Composición del corpus por etapa del pipeline. Elaboración propia.}
\label{tab:corpus}
\centering
\fontsize{9}{11}\selectfont
\setlength{\tabcolsep}{4pt}
\begin{tabular}{@{}p{2.6cm}rrp{1.6cm}@{}}
\toprule
\textbf{Etapa} & \textbf{Imágenes} & \textbf{Personas} & \textbf{17 kpts (\%)} \\
\midrule
Pool crudo (post-stride) & 74.107 & --- & --- \\
Pseudo-etiquetado & 74.107 & 73.955 & 94,3 \\
Derivado a auditoría (baja confianza) & 2.396 & --- & --- \\
Fold 0 — entrenamiento & --- & --- & --- \\
Fold 0 — validación & --- & --- & --- \\
\bottomrule
\end{tabular}
\end{table}
\pendiente{Completar las celdas "---" restantes: conteo train/val por fold =
contar líneas de los manifiestos \texttt{folds/fold0\_\{train,val\}.txt} (no
calculado todavía, ver PENDIENTES.md \S4); personas/\% completos del
subconjunto derivado a auditoría = depende de que se corra la auditoría real.}

\subsection{Costo de anotación}
Reportar el costo de la auditoría —tasa de corrección por motivo, cantidad de
cuadros efectivamente revisados y horas humanas insumidas— requiere haber
ejecutado la auditoría real sobre el subconjunto derivado, lo que todavía no
ocurrió. \pendiente{La auditoría CVAT no se ejecutó: no hay tasa de
corrección, cantidad de cuadros efectivamente auditados ni horas humanas
insumidas. Ver PENDIENTES.md \S6 para el comando exacto que generaría estos
números.}

\subsection{Configuración de entrenamiento y compresión}
La Tabla~\ref{tab:config} inventaría los hiperparámetros reales de cada fase.

\begin{table}[t]
\caption{Configuración de entrenamiento y compresión por fase. Elaboración propia.}
\label{tab:config}
\centering
\fontsize{9}{11}\selectfont
\setlength{\tabcolsep}{4pt}
\begin{tabular}{@{}p{2.0cm}p{2.6cm}p{1.6cm}@{}}
\toprule
\textbf{Fase} & \textbf{Parámetro} & \textbf{Valor} \\
\midrule
1 — Base & batch / optimizador & 48 / auto \\
2 — Poda & sparsity / fine-tune & 0,3 / 20 ép. \\
3 — Cuant. & calibración & 100--300 img. \\
4 — Destilación & $T$ / $\alpha$ / épocas & 4,0 / 0,5 / 100 \\
5 — Combinada (CV) & fine-tune / destil. & 10 ép. / 80 ép. \\
5 — Combinada (deploy) & fine-tune / destil. & 25 ép. / 20 ép. \\
\bottomrule
\end{tabular}
\end{table}

La Tabla~\ref{tab:cadena} reporta el resultado medido de aplicar esa cadena,
en dos corridas distintas que conviene no confundir entre sí: una corrida de
la fase 2 aislada (base, podado, recuperado) y la corrida de despliegue de la
fase 5, que retoma el modelo recuperado y continúa con la destilación.

\begin{table}[t]
\caption{Resultado medido de la cadena de compresión. Elaboración propia.}
\label{tab:cadena}
\centering
\fontsize{9}{11}\selectfont
\setlength{\tabcolsep}{3.5pt}
\begin{tabular}{@{}p{1.5cm}p{2.0cm}rrr@{}}
\toprule
\textbf{Corrida} & \textbf{Etapa} & \textbf{OKS} & \textbf{Params (M)} & \textbf{Lat. (ms)} \\
\midrule
Fase 2 aislada & Base (COCO) & 0,846 & 3,30 & 46,7 \\
 & Podado (sparsity 0,3) & 0,000 & 1,91 & 31,8 \\
 & Recuperado (ft 20 ép.) & 0,935 & 1,90 & --- \\
\addlinespace
Fase 5 deploy & Recuperado (ft 25 ép.) & 0,935 & --- & --- \\
 & Destilado (20 ép., roto) & 0,626 & --- & --- \\
\bottomrule
\end{tabular}
\end{table}

El estudiante recuperado supera al modelo maestro sobre la misma partición
(OKS 0,911 sobre una muestra de 600 cuadros), lo que confirma que la poda es
recuperable con el fine-tune adoptado; la fila de destilación de la fase 5
deploy, en cambio, es el punto de partida del diagnóstico de la subsección
siguiente y no debe leerse como el desempeño esperable de esa técnica.

\subsection{Diagnóstico de iteración negativa}
La primera iteración completa del pipeline de despliegue produjo dos
hallazgos que conviene reportar con el mismo detalle que un resultado
positivo, porque explican decisiones de diseño posteriores y porque un
hallazgo negativo bien diagnosticado sostiene la motivación de todo el
trabajo.

\emph{Colapso fuera de dominio.} El primer checkpoint desplegable (\emph{v2})
alcanzó un OKS de 0,931 sobre la partición de validación completa —un
resultado saludable— pero, al ejecutarse sobre video real, producía un
esqueleto visiblemente aplastado en el eje horizontal. La causa no fue la
cuantización: el modelo cuantizado, el modelo en punto flotante y el
checkpoint sin exportar coincidieron entre sí. Tampoco fue el encuadre:
cuadros completos del corpus de origen y recortes reencuadrados como los
produce la aplicación se comportaron con normalidad. Al evaluar el modelo
sobre 79 imágenes reales de un corpus general de pose —nunca vistas por el
estudiante— y normalizar la extensión horizontal de los puntos predichos por
el ancho de la caja de referencia, la razón cayó a 0,701 contra un techo de
0,977 alcanzado por un modelo entrenado solo sobre ese corpus general, con 40
de 76 imágenes efectivamente colapsadas y un error medio de posición de 0,314
de la diagonal de la caja de referencia: el esqueleto caía fuera del cuerpo,
no simplemente angosto. El diagnóstico señaló una falla de generalización de
apariencia del ajuste fino posterior a la poda, entrenado casi
exclusivamente sobre las imágenes homogéneas del corpus de rehabilitación. La
corrección consistió en repetir ese ajuste fino sobre una mezcla de ese mismo
corpus con una muestra del corpus general de pose, con aumentación de datos
reforzada; el checkpoint resultante (\emph{v3}) llevó la razón de 0,701 a
0,913, las imágenes colapsadas de 40/76 a 1/79, y de paso mejoró el OKS
in-domain de 0,931 a 0,941.

\emph{Destilación no resuelta.} La misma corrida de despliegue, al continuar
con la etapa de destilación tras el ajuste fino de recuperación, degradó al
estudiante de un OKS de 0,935 a 0,626 en veinte épocas, con predicciones de
coordenadas de keypoint sin valor numérico definido. El diagnóstico
identificó tres defectos concretos y no uno solo: la tasa de aprendizaje de
la destilación es diez veces la del ajuste fino que la precede, porque no se
reduce al reiniciar el entrenamiento sobre un modelo ya convergido; el
término de transferencia de atención se normaliza dividiendo por el área
espacial del mapa de activación, lo que lo reduce a menos del 0,5\,\% de la
pérdida total y lo vuelve numéricamente inerte; y un factor de temperatura al
cuadrado, heredado de la destilación clásica sobre distribuciones de
probabilidad \cite{hinton2015}, se aplica sin justificación a mapas de
atención que no son distribuciones. Los dos últimos defectos se cancelaban
entre sí, razón por la cual el problema no era evidente antes de instrumentar
cada término por separado: corregir la normalización sin corregir también el
factor de temperatura invertiría la falla, haciendo que el término de
destilación domine sobre la tarea principal. \pendiente{No usar destilación
de las fases 4 y 5 para resultados reales hasta corregir conjuntamente los
tres defectos y recalibrar el peso relativo de ambos términos de la
pérdida.}

% ======================= 7. LIMITACIONES ================================
\section{Limitaciones}
Los alcances de lo reportado deben enunciarse con precisión. La auditoría
manual está implementada y su mecanismo de control de calidad —el control de
ida y vuelta— está verificado, pero no se ejecutó todavía sobre el corpus
real derivado a revisión humana: ninguna cifra de este artículo sobre tasa de
corrección o costo de anotación existe todavía. La etapa de destilación de
conocimiento de las fases 4 y 5 tiene un defecto de implementación no
resuelto y no debe interpretarse como una medición válida de lo que esa
técnica puede aportar sobre este dataset. La restricción de uso que motiva la
anotación autoalojada se asume por diseño; los términos exactos de la
licencia del dataset UCO Physical Rehabilitation no están verificados contra
su fuente oficial, y ninguna afirmación de este artículo debe leerse como una
interpretación legal confirmada. Finalmente, el modelo que efectivamente se
despliega se entrenó una única vez sobre una partición fija: no existe
todavía una corrida de validación cruzada agrupada por sujeto sobre esa
configuración exacta —con la mezcla de corpus general de pose y la
aumentación reforzada del lineage corregido—, omitida por el costo de
cómputo que impone repetir el entrenamiento cinco veces.

% ======================= 8. CONCLUSIONES =================================
\section{Conclusiones y trabajo futuro}
Este trabajo describió la construcción del conjunto de datos y el
entrenamiento de un estimador de pose comprimido para ejercicios de
rehabilitación domiciliaria, con foco en dos decisiones que lo hacen viable
en un dominio sin datasets públicos adecuados y con restricciones de uso
sobre los datos disponibles: una estrategia de anotación en dos niveles que
concentra el esfuerzo humano en la cola de casos difíciles en lugar de
distribuirlo sobre todo el corpus, y una cadena de compresión cuyo orden y
cuya elección de técnica de recuperación —destilación en lugar de ajuste fino
supervisado— responden directamente a la limitación de anotación del dataset
de origen. La distinción entre el modelo desplegable y el instrumento de
medición, y el diagnóstico honesto de dos hallazgos negativos de la primera
iteración, completan una descripción del proceso que prioriza la
trazabilidad de cada cifra por sobre presentar un resultado más prolijo de lo
que el estado actual del trabajo permite.

El trabajo futuro sigue dos líneas concretas. La primera es completar la
validación cruzada agrupada por sujeto sobre la configuración exacta del
modelo finalmente desplegado, pospuesta hasta ahora por su costo de cómputo.
La segunda es la finalización de la campaña de auditoría manual actualmente
en curso, cuyo resultado —tasa de corrección, costo de anotación y un
conjunto dorado de evaluación— es el insumo que falta para cerrar el
diagnóstico de calidad del pseudo-etiquetado que este artículo deja
planteado.

% ======================= REFERENCIAS ===================================
\begingroup
\fontsize{9}{11}\selectfont
\bibliographystyle{unsrt}
\bibliography{referencias}
\endgroup

\end{document}
